M112 & D081 & mixed & carrier, evaluator, project custody, training-data independence, human independence unresolved \\
\bottomrule
\end{tabularx}
\end{table}

\section{The inherited substrate and the failure vocabulary}
\label{sec:mechanism}
The experiments deliberately use a small symbolic substrate because the objective is causal resolution rather than benchmark scale. A state contains a component registry, an interface width, an operator or candidate space, and content-addressed rules acquired by earlier stages. Demands are evaluated by a fixed host and a frozen three-feature observation vocabulary:
\begin{enumerate}
  \item $g_0$: the demand needs an unread signal;
  \item $g_1$: candidate search is exhausted for this demand;
  \item $g_2$: progress remains available on the operator axis.
\end{enumerate}
A feature row is one of the eight Boolean triples. An attribution program maps a row to a component to modify. The baseline \mzero is hardwired to the operator table. Generation one may add a rule that selects the signal interface. Generation two may add a rule that selects the candidate space. M111 adds a diagnostic policy that can intercept an attribution when the historical record says the row is underdetermined.

The lower Boolean interpreter initially holds only monotone operators \code{AND} and \code{OR}. Its complete image over two signals contains 4 of the 16 Boolean functions. The key structural tool is a monotonicity lemma: any expression composed solely from monotone operators remains monotone regardless of node budget. This lets the line distinguish an expressivity boundary from an insufficient search budget. Several fresh controls still run at larger node bounds, but the lemma is what licenses the stronger claim when a non-monotone target is outside the held language.

This substrate is intentionally unlike an LLM agent. It is deterministic, finite, and exhaustively inspectable. That limits external validity but enables forms of causal evidence that are difficult to obtain in large learned systems: complete function images, exact state equality, exhaustive attribution-row censuses, and semantic ablations whose consequences can be traced without stochastic confounds.

\section{M107: extending the lower interpreter}
\label{sec:m107}
M107 asks whether the lineage can change what its interpreter can express rather than merely select a different expression inside a fixed language. With \code{AND} and \code{OR}, the complete image is 4 of 16 Boolean functions. Two frozen demands are excluded by the monotonicity lemma at every node bound. One demand alone leaves two surviving reach classes and is refused. Joint demands exhaust a twenty-operator candidate space and identify one reach class. The lineage adopts a unary non-monotone operator with truth table $(1,0)$ as content-addressed state.

After adoption, the complete image becomes 16 of 16; both targets construct and execute after producer death. Ablation returns the image exactly to 4, and fresh controls still fail at bounds 9 and 13. The canonical checker computes P1--P16 true and replay equality, with zero model, network, or remote-execution calls.

The scientific point is not that negation is a novel operator. It is that the lower interpreter begins with an explicitly enumerated expressive ceiling and ends with a serialized operator whose registration changes that ceiling. The acquired definition is later recovered as predecessor state rather than rewritten into successor code. M108 and M111 rely on that distinction: later expressibility is causally linked to what the lineage previously acquired.

M107 also fixes the first rung of the terminology used throughout the paper. It is an \emph{interpreter extension}, not yet an acquisition-machinery modification. The machinery that observes a failure and decides which system component to change is still authored and hardwired. Thus M107 establishes a prerequisite for deeper adaptation without by itself demonstrating machinery that improves machinery.

\section{M108: modifying the machinery that performs later acquisitions}
\label{sec:m108}
M108 reconstructs M107's acquired operator rather than hand-coding a stronger baseline. At the base interface width, the operator table is already saturated at 16 of 16 Boolean functions. A later failure therefore cannot honestly be explained by a missing Boolean operator. The acquisition machinery nevertheless begins hardwired to blame that axis.

The lineage receives a bounded episode record over attribution rows. With episodes covering only rows $\{0,3\}$, two attribution classes survive and acquisition is refused. With rows $\{0,2,3\}$, the sixteen-program attribution space yields two consistent programs in one attribution class. The adopted rule has truth table $[1,0,1,0]$, is non-monotone, two nodes long, and becomes content-addressed lineage state.

On the held demand $B$, the modified state \mone attributes the failure to the signal interface, widens it, and constructs the target in six nodes. Fresh \mzero exhausts all twenty operators and refuses at both bounds 9 and 13 under an equal machinery-step budget. Byte-exact ablation returns to \mzero and removes the capability; mutation sends attribution back to the operator axis; corruption fails closed on the state identity mismatch. P1--P16 compute true, replay is equal, and the canonical run uses fourteen isolated processes with zero external calls.

Two structural arguments make the dependency stronger than ``the search happened not to find a solution.'' First, at the base interface width, all 16 liftable images of the 256 world functions exclude $B$ for any operator set. Second, at full width, $B$ lies outside the monotone image by lemma. The capability therefore requires both the earlier interpreter extension and the acquired attribution rule that tells the machinery to widen the interface.

